%% CHAPITRE 6 -- Tests, Observabilite et Resultats

\section*{Introduction}

Ce chapitre clôture la description du travail réalisé en présentant les résultats mesurables
de la transformation. Nous y décrivons la stratégie de tests multi-niveaux (unitaires, sécurité
et E2E), l'architecture d'observabilité avec Prometheus et Grafana, la démonstration des
vulnérabilités intentionnelles, et enfin le bilan comparatif avant/après avec les indicateurs
de performance clés (KPIs) qui quantifient la valeur ajoutée de la plateforme DevSecOps.

\section{Strategie de Tests Multi-Niveaux}
La pyramide de tests du projet comprend trois niveaux complementaires~:
\begin{enumerate}
  \item \textbf{Tests unitaires} (Vitest + Testing Library)~: tests rapides, nombreux,
        isolent les composants React individuellement~;
  \item \textbf{Tests de securite} (Trivy, ZAP, SonarCloud)~: tests automatises de securite
        integres dans le pipeline CI/CD~;
  \item \textbf{Tests E2E} (Playwright in-cluster)~: tests lents mais a haute valeur ajoutee,
        valident le comportement reel de l'application dans le contexte de production.
\end{enumerate}
\section{Tests Unitaires avec Vitest}
Les tests unitaires couvrent les composants React critiques~:
\begin{itemize}
  \item \textbf{TodoItem}~: rendu du texte, toggle completed, suppression~;
  \item \textbf{AuthContext}~: login/logout, persistence de session~;
  \item \textbf{TodoDashboard}~: CRUD complet, filtrage, compteurs.
\end{itemize}
La couverture de code est superieure a 80\% grace au Quality Gate SonarCloud qui bloque
toute PR faisant baisser la couverture sous ce seuil.
\section{Tests E2E Playwright --- 6 Scenarios}
\begin{table}[h]
\centering
\setlength{\arrayrulewidth}{0.8pt}
\setlength{\dashlinedash}{3pt}
\setlength{\dashlinegap}{2pt}
\renewcommand{\arraystretch}{1.4}
\caption{Les 6 scenarios de tests Playwright E2E}
\label{tab:playwright}
\begin{tabularx}{\textwidth}{|>{\centering\arraybackslash}X|>{\centering\arraybackslash}X|>{\centering\arraybackslash}X|}
\hline
\tableheader \textbf{Tag} & \textbf{Scenario} & \textbf{Ce que ca valide} \\
\hline
@t1 & Credentials invalides -> message d'erreur & Authentification, gestion erreurs \\ \hdashline
@t2 & Login -> dashboard vide affiche & Flux d'authentification complet \\ \hdashline
@t3 & Login -> Logout -> retour formulaire & Cycle auth complet \\ \hdashline
@t4 & Session persiste apres rechargement & Persistance localStorage \\ \hdashline
@t5 & Todos persistent apres rechargement & Persistance des donnees \\ \hdashline
@t6 & Ajouter, completer, supprimer une tache & CRUD complet \\
\hline
\end{tabularx}
\end{table}
Les tests sont executes en trois modes~:
\begin{itemize}
  \item \textbf{In-cluster (PostSync)}~: jobs Kubernetes crees par Argo CD apres chaque
        deploiement~;
  \item \textbf{Local CI (matrice)}~: GitHub Actions execute Playwright avec \texttt{--grep
        @t1..@t6}~;
  \item \textbf{Standalone Job}~: \texttt{k8s/testing/playwright-job.yaml} pour tests manuels
        ponctuels.
\end{itemize}
\section{Observabilite avec Prometheus et Grafana}
\subsection{Architecture d'Observabilite}
Le flux de collecte des metriques~:
\begin{enumerate}
  \item Le sidecar \texttt{nginx-exporter} (port 9113) expose les metriques Nginx~;
  \item Le \texttt{ServiceMonitor} dans le namespace \texttt{observability} scrape ces metriques
        toutes les 30 secondes~;
  \item Prometheus stocke les series temporelles~;
  \item Grafana interroge Prometheus pour les tableaux de bord.
\end{enumerate}
\subsection{Metriques de Securite Cles}
\begin{table}[h]
\centering
\setlength{\arrayrulewidth}{0.8pt}
\setlength{\dashlinedash}{3pt}
\setlength{\dashlinegap}{2pt}
\renewcommand{\arraystretch}{1.4}
\caption{Metriques Nginx collectees et leur utilite securite}
\label{tab:metrics}
\begin{tabularx}{\textwidth}{|>{\centering\arraybackslash}X|>{\centering\arraybackslash}X|>{\centering\arraybackslash}X|}
\hline
\tableheader \textbf{Metrique} & \textbf{Description} & \textbf{Utilite Securite} \\
\hline
nginx\_connections\_active & Connexions actives & Detection de DDoS \\ \hdashline
nginx\_http\_requests\_total & Requetes par status code & Spike de 4xx = attaque de scan \\ \hdashline
nginx\_connections\_accepted & Connexions acceptees total & Baseline trafic normal \\ \hdashline
nginx\_connections\_handled & Connexions traitees & Performance et disponibilite \\
\hline
\end{tabularx}
\end{table}
\section{Resultats et Comparaison Avant/Apres}
\subsection{Tableau de Bord de Securite}
\begin{table}[h]
\centering
\setlength{\arrayrulewidth}{0.8pt}
\setlength{\dashlinedash}{3pt}
\setlength{\dashlinegap}{2pt}
\renewcommand{\arraystretch}{1.4}
\caption{Bilan securite avant et apres la transformation DevSecOps}
\label{tab:bilan_securite}
\begin{tabularx}{\textwidth}{|>{\centering\arraybackslash}X|>{\centering\arraybackslash}X|>{\centering\arraybackslash}X|}
\hline
\tableheader \textbf{Domaine} & \textbf{Avant (Jenkins)} & \textbf{Apres (DevSecOps)} \\
\hline
SAST & Absent & SonarCloud detecte XSS, hotspots \\ \hdashline
SCA & Absent & Trivy FS : 9 CVEs sur 3 packages \\ \hdashline
Container Scan & Absent & Trivy Image : 47 CVEs nginx:1.21 \\ \hdashline
IaC Scan & Absent & Trivy Config : 8 misconfigs K8s \\ \hdashline
DAST & Absent & OWASP ZAP : 10 findings securite \\ \hdashline
Secrets Scan & Absent & Trivy : secrets ENV + .env detectes \\ \hdashline
Runtime Security & Absent & Wazuh : 4 agents actifs \\ \hdashline
Reseau & Tout ouvert & Zero Trust : 12 NetworkPolicies \\ \hdashline
Auth CI & Cles JSON statiques & WIF OIDC : zero cle persistante \\ \hdashline
Tests & Unit tests Jenkins & Unit + E2E + 6 scenarios Playwright \\ \hdashline
Infrastructure & ClickOps & Terraform IaC : 100\% reproductible \\ \hdashline
Deploiement & Scripts manuels & Argo CD GitOps : auto-sync + self-heal \\ \hdashline
Audit trail & Inexistant & Chaque commit/build/deploy trace \\
\hline
\end{tabularx}
\end{table}
\subsection{KPIs Techniques Mesures}
\begin{table}[h]
\centering
\setlength{\arrayrulewidth}{0.8pt}
\setlength{\dashlinedash}{3pt}
\setlength{\dashlinegap}{2pt}
\renewcommand{\arraystretch}{1.4}
\caption{Indicateurs cles de performance avant/apres}
\label{tab:kpis}
\begin{tabularx}{\textwidth}{|>{\centering\arraybackslash}X|>{\centering\arraybackslash}X|>{\centering\arraybackslash}X|}
\hline
\tableheader \textbf{KPI} & \textbf{Avant} & \textbf{Apres} \\
\hline
Vulnerabilites detectees avant prod & 0 & 9+ CVEs + 10 findings ZAP \\ \hdashline
Temps de detection d'une CVE critique & Jamais & < 5 min (sur push) \\ \hdashline
Infrastructure reproductible & Non & Oui (terraform apply) \\ \hdashline
Rollback automatique & Non & Oui (git revert + Argo CD) \\ \hdashline
Tests E2E automatises & 0 & 6 scenarios in-cluster \\ \hdashline
Couverture tests unitaires & 40\% & > 80\% (gate SonarCloud) \\ \hdashline
Auth CI sans cle statique & Non & Oui (WIF OIDC) \\ \hdashline
Namespaces avec isolation reseau & 0 & 5 (100\% du cluster) \\ \hdashline
SIEM actif & Non & Oui (Wazuh 4 agents) \\
\hline
\end{tabularx}
\end{table}
\section{Limites et Plan d'Amelioration Continue}
\subsection{Limites Actuelles}
\begin{itemize}
  \item \textbf{GKE Autopilot et Wazuh}~: hostPID/hostNetwork bloques, monitoring kernel
        limite~;
  \item \textbf{Authentification applicative}~: credentials hardcodes cote client (demo)~;
  \item \textbf{Secrets applicatifs}~: dans les variables ENV (demo intentionnelle)~;
  \item \textbf{Single environment}~: un seul environnement (production)~;
  \item \textbf{DAST baseline}~: scan passif uniquement, pas de scan actif complet.
\end{itemize}
\subsection{Roadmap Court Terme (1--3 mois)}
\begin{itemize}
  \item Ajouter CodeQL pour SAST JavaScript avance~;
  \item Implementer cosign pour la signature des images (SLSA Level 2)~;
  \item Generer un SBOM pour chaque image (syft ou trivy sbom)~;
  \item Ajouter Gitleaks en pre-commit pour bloquer les secrets avant push~;
  \item Ajouter Kyverno pour le Policy-as-Code en admission control.
\end{itemize}
\subsection{Roadmap Moyen Terme (3--6 mois)}
\begin{itemize}
  \item Migrer vers GCP Secret Manager + External Secrets Operator~;
  \item Ajouter un environnement staging avec promotion d'artefacts~;
  \item Implementer Checkov pour le scan IaC Terraform~;
  \item Configurer kube-bench pour le score de conformite CIS Kubernetes~;
  \item Ajouter des alertes Grafana de securite (Prometheus Alertmanager).
\end{itemize}
\subsection{Alignement aux Frameworks de Securite}
\begin{table}[h]
\centering
\setlength{\arrayrulewidth}{0.8pt}
\setlength{\dashlinedash}{3pt}
\setlength{\dashlinegap}{2pt}
\renewcommand{\arraystretch}{1.4}
\caption{Alignement aux frameworks de securite}
\label{tab:frameworks}
\begin{tabularx}{\textwidth}{|>{\centering\arraybackslash}X|>{\centering\arraybackslash}X|>{\centering\arraybackslash}X|}
\hline
\tableheader \textbf{Framework} & \textbf{Couverture Actuelle} & \textbf{Objectif} \\
\hline
OWASP Top 10 & A03, A05, A06, A09 couverts & A01, A02, A08 a ajouter \\ \hdashline
NIST SSDF & PW.1 (Build), RV.1 (Scan) & PW.4 (Signing), PO.3 (Reqs) \\ \hdashline
CIS Controls v8 & Config, Vuln Mgmt, AppSec & Inventory, Logging \\ \hdashline
OWASP SAMM & Secure Build, Sec Testing & Secure Architecture, Operations \\
\hline
\end{tabularx}
\end{table}

\section*{Conclusion}

Ce chapitre a fourni la synthèse quantitative de la transformation~: six scénarios Playwright
validés en cluster, une couverture de tests unitaires supérieure à 70~\%, un tableau de bord
Grafana opérationnel, et un bilan avant/après qui démontre la suppression de toutes les lacunes
sécuritaires identifiées au chapitre~1. Ces résultats confirment que la problématique posée
en début de stage a été pleinement résolue. La conclusion générale du rapport revient sur
l'ensemble de la démarche et ouvre des perspectives d'amélioration continue.

